% ================================================================
% 03_theorem.tex — 定理环境与引用帧
% ================================================================

\begin{frame}{定理环境与引用}
  \begin{theorem}{欧拉恒等式}{thm:euler}
    对任意实数 $\theta$, 有 $e^{i\theta} = \cos\theta + i\sin\theta$。
  \end{theorem}

  \pause

  \begin{proof}
    由 $e^{i\theta}$ 的泰勒展开直接得到, 见 \ref{thm:euler} 的证明过程。
  \end{proof}
	$g^{\prime}\left(x\right)=2x\implies g\left(x\right)=x^{2}+C$,其中$C\in\R$.
\end{frame}

\begin{frame}{理想的定义}
	\begin{definition}{理想}{}
		设$A$是环,$\mathfrak{a}$是$A$的加法子群.
		\begin{enumerate}
			\item 如果$A\mathfrak{a}\subseteq\mathfrak{a}$,则称$\mathfrak{a}$是$A$的左理想.
			\item 如果$\mathfrak{a}A\subseteq\mathfrak{a}$,则称$\mathfrak{a}$是$A$的右理想.
			\item 如果$A\mathfrak{a}\subseteq\mathfrak{a},\mathfrak{a}A\subseteq\mathfrak{a}$,则称$\mathfrak{a}$是$A$的双边理想.
		\end{enumerate}
	\end{definition}
	熟知环$A$是典范的$\left(A,A\right)$-双模.照此观点,
	\begin{itemize}
		\item $A$的左理想是左$A$-模$A_{s}$的子模.
		\item $A$的右理想是右$A$-模$A_{d}$的子模.
		\item $A$的双边理想是$\left(A,A\right)$-双模$_{s}A_{d}$的子双模.
	\end{itemize}
\end{frame}
